\chapter{Proposed System}
\label{ch:system}

This chapter presents the proposed system, a \emph{tiered repair bot}.
Section~\ref{sec:sys-req} states what a practical repair bot needs.
Section~\ref{sec:sys-overview} shows the bot. The next four sections describe
its modules, and Section~\ref{sec:sys-regimes} states when it helps.

\section{Usage Scenario and Requirements}
\label{sec:sys-req}

\paragraph{Scenario.} A test fails in continuous integration. A repair bot is
given the faulty code and tries to fix it within a budget set by the team. It
returns either a plausible patch for a developer to review, or a report of why
it gave up. Developers name CI as the place they most want repair to
run~\citep{noller2022trust}, and deployed systems work this way: they act on
failures found in CI and send every fix to a
reviewer~\citep{kirbas2021bloomberg,maddila2025agentic,marginean2019sapfix}.
Because it runs on every failing build, the bot meets easy and hard bugs
alike, and it cannot know in advance which is which.
Table~\ref{tab:sys-needs} lists what such a bot needs.

\begin{table}[htbp]
\centering
\caption{What a repair bot in CI needs, and the evidence.}
\label{tab:sys-needs}
\footnotesize
\begin{tabularx}{\textwidth}{@{}>{\raggedright\arraybackslash}p{0.19\textwidth} >{\raggedright\arraybackslash}X >{\raggedright\arraybackslash}p{0.27\textwidth}@{}}
\toprule
Requirement & Evidence & Served by\\
\midrule
P1 bounded time & $63\,\%$ of $103$ practitioners would not wait over an hour~\citep{noller2022trust}, though correctness matters more than speed~\citep{winter2023feel} & M0 Warm Substrate, M1 Evidence Memory\\
P2 predictable, low cost & budgets decided the adoption of repair at Bloomberg~\citep{kirbas2021bloomberg} & M2 Escalation Controller, M3 Budget Guard\\
P3 drop-in & developers want repair inside the CI pipeline~\citep{noller2022trust}, fitting their workflow~\citep{winter2022bloomberg} & no training data; one coverage run of the faulty program\\
P4 trustworthy & a fix must pass all CI checks and be reviewed~\citep{kirbas2021bloomberg} & M0 and M1: every verdict comes from the tests\\
P5 explain a failure & unexplained fixes are ignored~\citep{winter2022bloomberg}; developers value the fault location and tests~\citep{noller2022trust} & M3 Budget Guard\\
\bottomrule
\end{tabularx}
\end{table}

The evidence is strong for P3 and P4 and mixed for P1, so the bot aims at a
bounded time, not the shortest. P2 and P5 are supported only indirectly; they
are design choices the evidence favours.

\section{System Overview}
\label{sec:sys-overview}

Figure~\ref{fig:sys} shows one session. Before the first round the
\textbf{Warm Substrate} (M0) starts the test runtime on the project as the
failing build left it; the runtime then stays alive for the whole session. The
cheap model proposes candidates. The \textbf{Evidence Memory} (M1) checks each
one: a candidate seen before gets its stored verdict, and a new one is
compiled by the substrate, loaded into the live runtime and run against the
scheduled tests until the first failure. After a rejection, the \textbf{Budget
Guard} (M3) checks that budget is left, and the \textbf{Escalation Controller}
(M2) checks whether the candidate made progress. While the cheap tier makes
progress, it proposes again. When it stalls, the controller writes a handover
note and the strong model takes over. If the strong tier stalls too, or the
budget is spent, the session ends with a give-up report.

The repair loop interacts with the bot through a single call. It takes a
candidate and returns the verdict, the evidence, and what the session should
do next:
\begin{Verbatim}[fontsize=\footnotesize,xleftmargin=8mm]
result = bot.validate(candidate)
result.verdict    # plausible, rejected, does not compile, or repeat
result.evidence   # the failing test and its failure signature
result.next       # same tier, step up, or stop
result.note       # the handover note, when next is step up
result.report     # the give-up report, when next is stop
\end{Verbatim}
The bot is configured with a list of tiers and a budget per bug. It needs no
training data.

\begin{figure}[htbp]
\centering
\resizebox{\textwidth}{!}{%
\begin{tikzpicture}[
  font=\sffamily\small,
  box/.style={draw=boxline,line width=0.6pt,rounded corners=2.5pt,align=center,inner sep=4pt,
              minimum height=12mm,text width=38mm,fill=white},
  hi/.style={box,draw=hiline,line width=1pt,fill=hifill},
  tier/.style={box,text width=40mm},
  term/.style={box,text width=20mm},
  lab/.style={font=\sffamily\scriptsize,text=black!70,align=left},
  arr/.style={-{Stealth[length=2.2mm]},draw=black!45,line width=0.7pt}
]
\node[font=\sffamily\small\bfseries,anchor=west] at (-2.6,2.75)
  {Tests check every candidate; a verified stall moves the session up a tier};
% the two tiers, inside one container
\draw[draw=boxline,line width=0.6pt,rounded corners=3pt,fill=black!6] (-2.6,-1.95) rectangle (2.6,1.95);
\node[font=\sffamily\small\bfseries,anchor=north west] at (-2.45,1.9) {Model tiers};
\node (t1) at (0,0.55)   [tier] {\textbf{Tier 1: cheap model}\\[-1pt]{\scriptsize\color{black!65} $n$ blind samples per call}};
\node (t2) at (0,-1.05)  [tier] {\textbf{Tier 2: strong model}\\[-1pt]{\scriptsize\color{black!65} conversation, feedback}};
% the modules
\node (m1) at (6.6,0)     [box,text width=42mm] {\textbf{M1 Evidence Memory}\\[-1pt]{\scriptsize\color{black!65} verdicts, stored tests}};
\node (m0) at (12.2,0)    [box,text width=42mm] {\textbf{M0 Warm Substrate}\\[-1pt]{\scriptsize\color{black!65} compile, load, run tests}};
\node (pl) at (17.0,0)    [term] {\textbf{Plausible}\\ \textbf{patch}};
\node (m3) at (12.2,-3.7) [box,text width=42mm] {\textbf{M3 Budget Guard}\\[-1pt]{\scriptsize\color{black!65} budget left, not stuck?}};
\node (gu) at (17.0,-3.7) [term] {\textbf{Give-up}\\ \textbf{report}};
\node (m2) at (5.9,-3.7)  [hi,text width=52mm] {\textbf{M2 Escalation Controller}\\[-1pt]{\scriptsize\color{black!65} progress, or a stall?}};
% arrows
\draw[arr] (2.6,0) -- node[above=1pt,lab]{candidate} (m1.west);
\draw[arr] (m1) -- node[above=1pt,lab]{new} (m0);
\draw[arr] (m0) -- node[above=1pt,lab]{all pass} (pl);
\draw[arr] (m0) -- node[right=2pt,lab,pos=0.5]{rejected} (m3);
\coordinate (rj) at ([xshift=-12mm]m3.north);
\draw[arr] (m1.south) -- (m1.south |- 0,-1.9) -- node[above=1pt,lab,pos=0.5]{repeat: stored verdict} (rj |- 0,-1.9) -- (rj);
\draw[arr] (m3) -- node[above=1pt,lab]{no} (gu);
\draw[arr] (m3) -- node[above=1pt,lab]{yes} (m2);
\draw[arr] (m2.west) -- (0,-3.7) -- (0,-1.95);
\node[lab,anchor=north east,align=right] at (-0.2,-2.55) {progress: same tier\\[1pt] stall: step up, with a note};
\end{tikzpicture}}
\caption[One repair session with the tiered bot.]{One repair session with the
tiered bot. Every candidate of either tier passes through the Evidence Memory
and the Warm Substrate. The Escalation Controller (highlighted) is the new
module: while a tier makes progress it proposes again, and a verified stall
moves the session up a tier with a handover note.}
\label{fig:sys}
\end{figure}

\begin{table}[htbp]
\centering
\caption[The four modules.]{The four modules: what each improves, the cost it
targets, and how.}
\label{tab:sys-modules}
\footnotesize
\begin{tabularx}{\textwidth}{@{}>{\raggedright\arraybackslash}p{0.15\textwidth} >{\raggedright\arraybackslash}p{0.24\textwidth} >{\raggedright\arraybackslash}p{0.23\textwidth} >{\raggedright\arraybackslash}X@{}}
\toprule
Module & Improves & Targets & How\\
\midrule
M0 Warm Substrate & the build and the start-up are paid once per session, not once per candidate & the build and start-up parts of validation, and the tests that do not reach the change & in-process compilation of the patched file; a fresh class loader in a live runtime; a fixed test schedule\\
M1 Evidence Memory & nothing is validated twice, and evidence survives a model switch & the validation of repeats and of candidates a stored test refutes & normalised verdict cache; stored failing tests and failure signatures\\
M2 Escalation Controller & the strong model is called only when it is needed & strong-tier calls on bugs the cheap tier repairs & cheap tier first; \Step\ moves up on a verified stall; handover note\\
M3 Budget Guard & predictable cost, and a reason when it fails & sessions that are never repaired & a budget per tier and per session; a stop when the top tier stalls; give-up report\\
\bottomrule
\end{tabularx}
\end{table}

\paragraph{A session, step by step.} Table~\ref{tab:sys-walk} follows one bug
through the bot.

\begin{table}[htbp]
\centering
\caption{One session of a bug in \texttt{DateUtils} (illustrative).}
\label{tab:sys-walk}
\footnotesize
\begin{tabularx}{\textwidth}{@{}c >{\raggedright\arraybackslash}p{0.08\textwidth} >{\raggedright\arraybackslash}p{0.16\textwidth} >{\raggedright\arraybackslash}X >{\raggedright\arraybackslash}p{0.17\textwidth}@{}}
\toprule
Step & Tier & Candidates & Evidence Memory & Controller\\
\midrule
1 & cheap & five from one call & two are the same patch after normalising; of the four distinct ones, one does not compile, two fail \texttt{DateUtilsTest} with error A and one fails \texttt{CalendarTest} with error B & progress: three new signatures\\
2 & cheap & five more & one repeats a patch of step 1 and is answered from the memory; four are new, and all fail with error A or B & no new signature: stall, step up\\
3 & strong & one, given the handover note & fails \texttt{DateUtilsTest} with a new error C & progress\\
4 & strong & one, given the failure & passes every test & plausible: stop\\
\bottomrule
\end{tabularx}
\end{table}

The strong model was called twice. The ten cheap candidates cost two calls of
the cheap model and eight validations on the warm substrate; two of them were
answered from the memory.

\section{Module M0: Warm Validation Substrate}
\label{sec:sys-warm}

\paragraph{Principle.} During a repair session the project, its dependencies
and the test runtime never change. One file changes, in the bodies of a few
methods. A validator that rebuilds the project and starts a fresh runtime for
every candidate pays, once per candidate, for work whose result is the same
every time. A cheap-first bot multiplies candidates, so it pays this more
often than any other design.

\paragraph{Mechanism.} M0 does this work once per session and keeps it.
\begin{enumerate}
\item \textbf{Build once.} The faulty program is compiled once, with the
project's own build. In CI the failing build has already done this.
\item \textbf{Compile one file, in process.} For each candidate only the
patched source file is compiled, by the compiler shipped with the runtime
(\code{javax.tools}), inside a process that stays alive, against the classes
of the first build.
\item \textbf{Load, do not restart.} The new bytecode is loaded into the live
test runtime through a fresh class loader, and the tests are run there. The
Java runtime and the test framework stay loaded and warm; the classes of the
project, its tests and its libraries are loaded again for each candidate.
\end{enumerate}

\paragraph{Why a fresh class loader.} On-the-fly validators for classical
repair redefine the patched class in place~\citep{chen2021uniapr}.
Redefinition has two limits in this setting. It cannot change the shape of a
class, and a candidate that only rewrites a method body changes it as soon as
the new body holds a lambda or an anonymous class. And it keeps the static
state of every loaded class, so what one candidate leaves behind has to be
reset by other means: UniAPR re-runs class initialisers between patches, and
before it did so it disagreed with ordinary validation on $4$ of $74$
bugs~\citep{chen2021uniapr}. A class loader per candidate gives both
properties by construction. Any candidate that compiles can be loaded, and no
static state of the project, its tests or its libraries survives from one
candidate to the next. The price is that these classes are loaded again each
time; the build and the start-up are still paid once.

\paragraph{Calibration.} At the start of a session the substrate compiles the
unchanged file in the live runtime and runs the scheduled tests once. Every
test must end as it did in the failing CI run: the tests that failed there
fail, and the others pass. If one differs, the live runtime is not trusted for
this bug, for instance because a test depends on how its classes are loaded,
and the session falls back to the cold path. The run also warms the runtime up
before the first candidate arrives.

\paragraph{Test schedule.} The substrate runs the tests in a fixed order that
needs no learning: first the tests that failed on the faulty program, then the
tests that rejected earlier candidates of the session, then the remaining
tests that cover the faulty methods, fastest first. Tests that do not cover
the faulty methods are skipped, which is safe when a patch only replaces the
bodies of those methods. The first step is common practice, the second is
FRTP~\citep{qi2013frtp}, and the third is the rule of
PIT~\citep{coles2016pit}. What this schedule leaves to any other order is
measured (Section~\ref{sec:exp-e1}).

\paragraph{Reported failure.} A candidate can fail several tests, and the
order decides which failure is met first. What is reported about a rejected
candidate, to the model and in its failure signature, is therefore fixed by
the static part of the order: its first failing test among those that failed
on the faulty program, and otherwise among the covering tests, fastest first.
The tests that rejected earlier candidates change what a rejection costs,
never what is reported. A candidate is thus told the same thing whatever was
tried before it, and this is what allows a session to be replayed
(Section~\ref{sec:exp-e2}). A deployed bot, which stops at the first failure
it meets, can differ from this rule only on a candidate that passes every
originally failing test and then fails several others; how often that happens
is measured.

\paragraph{Safety.} M0 must not change any verdict. Three things could.
\begin{itemize}
\item \textbf{State that outlives a class loader}: system properties, the
default locale and time zone, files, threads. The substrate restores the first
three after each candidate and interrupts the threads a candidate leaves
behind.
\item \textbf{A candidate that hangs or kills the runtime.} A test that passes
its time limit cannot be stopped inside a JVM, and a candidate may exit the
process. The substrate counts that test as failed, discards the runtime and
starts a new one, and the new start is charged to the candidate.
\item \textbf{Tests that depend on how classes are loaded.} The calibration
run catches them before the first candidate.
\end{itemize}
The question is treated as empirical. Every candidate of the study is
validated on the cold path as well, the verdict of record is the cold one, and
every disagreement is reported (Section~\ref{sec:exp-measure}).

\paragraph{Expected size.} Table~\ref{tab:sys-m0} sets the two paths side by
side. UniAPR and ExpressAPR report speed-ups of one to two orders of magnitude
on classical repair tools~\citep{chen2021uniapr,xiao2024expressapr}, where
most of what they remove is this same build and start-up. How much of that
carries over to an LLM loop, which validates tens of candidates per bug
rather than thousands, is what RQ2 measures. For escalation the consequence is
direct: on the cold path every extra cheap candidate pays a build and a
start-up of its own.

\begin{table}[htbp]
\centering
\caption{What one candidate pays with and without M0.}
\label{tab:sys-m0}
\footnotesize
\begin{tabularx}{\textwidth}{@{}>{\raggedright\arraybackslash}p{0.27\textwidth} >{\raggedright\arraybackslash}p{0.3\textwidth} >{\raggedright\arraybackslash}X@{}}
\toprule
Step & Without M0 (cold path) & With M0 (warm path)\\
\midrule
Build the candidate & the project's own build & one file, compiled in process\\
Start the test runtime & once per candidate & once per session, with the calibration run\\
Run the scheduled tests & to the first failure & the same tests, in a runtime that is already warm\\
Candidate that hangs or kills the runtime & the runtime is discarded anyway & a new runtime is started and charged to the candidate\\
\bottomrule
\end{tabularx}
\end{table}

\paragraph{Relation to existing work.} On-the-fly patch validation is not
new~\citep{chen2021uniapr,xiao2024expressapr}. M0 claims no new mechanism. Its
role here is to make the cheap tier's candidates cheap to check, and its
contribution is the measurement of what that is worth in an LLM loop (RQ2).

\section{Module M1: Evidence Memory}
\label{sec:sys-memory}

\paragraph{Principle.} What the tests have shown about a candidate is true
whichever model proposed it. The memory keeps it for the whole session and
across the switch between models.

\paragraph{What it keeps.}
\begin{itemize}
\item \textbf{Verdicts.} Each candidate is first \emph{normalised}: it is
parsed into a syntax tree, comments are removed, local variables are renamed
in order of appearance, and the tree is printed back. A candidate whose
normalised form is already in the memory is a \emph{repeat}: it gets the
stored verdict without running any test.
\item \textbf{Failing tests.} Every test that has rejected a candidate is kept
and scheduled first for later candidates.
\item \textbf{Failure signatures.} Each distinct signature
(Section~\ref{sec:bg-stall}) is kept with a count and one example patch. The
controller reads progress from them, and the handover note and the give-up
report are written from them.
\end{itemize}

\paragraph{Across tiers.} The memory outlives the switch. A strong-tier
candidate that equals a cheap-tier one is answered from the memory, and one
that fails a stored test is rejected by that test first. How often the strong
model repeats what the cheap one already showed to be wrong is itself a
measurement (RQ4). One rule keeps the tiers separable: a model is told when it
repeats \emph{itself}, but a repeat of the other tier's candidate is reported
to it as an ordinary failure. Under a cold handover the strong model is
therefore shown exactly what it would be shown if it ran alone.

\paragraph{Safety.} Normalisation removes only differences that cannot affect
behaviour: layout, comments and the names of local variables. Two patches
therefore share a verdict only when they behave identically, so reusing it is
safe for deterministic tests. Running stored failing tests first changes the
order of the tests, not the set, so it never changes a verdict.

\section{Module M2: Escalation Controller}
\label{sec:sys-controller}

\paragraph{Principle.} A bug should be repaired by the cheapest model that can
repair it. Which model that is cannot be known in advance, but it can be found
out cheaply: the cheap model tries first, and the tests say whether it is
getting anywhere.

\paragraph{Tier 1: blind re-sampling.} The cheap tier does not hold a
conversation. Each call sends the same first prompt and asks for $n$
candidates; nothing about earlier failures is added. Two results support this
choice. Independent re-prompting repairs more bugs than a conversation at the
same number of attempts~\citep{chen2025conversational}. And small models
anchor on their own failed attempt~\citep{verma2026blind}. Blind re-sampling
also keeps the prompt constant, so one call returns several candidates for the
price of one prompt~\citep{hidvegi2024cigar}.

\paragraph{Tier 2: conversation with feedback.} The strong tier runs a
conversational loop in the style of
ChatRepair~\citep{xia2024chatrepair}: after each rejected candidate,
the failing test and its error are fed back. What such feedback gains depends
on the model that reads it~\citep{olausson2024selfrepair}, which is why only
the strong tier receives it.

\paragraph{The rule.} After each cheap call, the controller counts the
candidates that made progress (Section~\ref{sec:bg-stall}). A call has
\emph{stalled} when at most one of its candidates made progress. The session
steps up after a stalled call, or when the cheap tier's budget of candidates
is spent. On a step-up the controller writes the handover note and starts the
strong tier. Algorithm~\ref{alg:step} gives the rule, called \Step. The
candidates of a call are validated in the order they were returned, and the
session ends at the first plausible one.

\begin{algorithm}[htbp]
\caption{\Step{} for one session}
\label{alg:step}
\begin{algorithmic}[1]
\Require cheap model $m_1$, strong model $m_2$; $n$ candidates per cheap call; cheap budget $B_1$ candidates; strong budget $R_2$ rounds
\State $\mathit{seen} \gets \emptyset$ \Comment{failure signatures, kept by the Evidence Memory}
\While{fewer than $B_1$ cheap candidates were drawn}
  \State draw $n$ candidates from $m_1$ with the first prompt
  \State $\mathit{new} \gets 0$
  \For{each candidate $p$}
    \State $v \gets \Call{Validate}{p}$ \Comment{the memory, then the schedule}
    \State \textbf{if} $v$ is plausible \textbf{then return} $p$
    \State \textbf{if} $\mathrm{signature}(v) \notin \mathit{seen}$ \textbf{then} add it to $\mathit{seen}$;\; $\mathit{new} \gets \mathit{new} + 1$
  \EndFor
  \State \textbf{if} $\mathit{new} \le 1$ \textbf{then break} \Comment{the call stalled}
\EndWhile
\State $\mathit{note} \gets \Call{Handover}{\mathit{seen}}$ \Comment{empty under a cold handover}
\For{$\mathit{round} = 1$ \textbf{to} $R_2$}
  \State $p \gets$ next candidate from $m_2$, given $\mathit{note}$ and the feedback so far
  \State $v \gets \Call{Validate}{p}$
  \State \textbf{if} $v$ is plausible \textbf{then return} $p$
  \State \textbf{if} the Budget Guard stops the session \textbf{then break} \Comment{Section~\ref{sec:sys-guard}}
\EndFor
\State \textbf{return} the give-up report
\end{algorithmic}
\end{algorithm}

\paragraph{The handover note.} The note is written by the bot from the
Evidence Memory, not by the cheap model. It holds the number of candidates
tried; each distinct failure signature, most frequent first, with its test,
its error and the number of candidates that ended there; and the compiler
errors seen. It holds none of the cheap model's explanations.
Figure~\ref{fig:sys-note} shows an example. Four handovers are compared (RQ4):
\begin{itemize}
\item \textbf{Cold}: no note. The strong tier starts from the first prompt
alone.
\item \textbf{Evidence}: the note described above.
\item \textbf{Evidence and patches}: the note, with the shortest candidate
that produced each signature.
\item \textbf{Transcript}: every cheap candidate and its failure, in the order
they occurred.
\end{itemize}

\begin{figure}[htbp]
\centering
\begin{minipage}{0.95\textwidth}
\begin{Verbatim}[frame=single,rulecolor=\color{boxline},fontsize=\footnotesize,framesep=2mm]
Handover note (10 candidates from the first model, none plausible)
Verified failures, most frequent first:
  6x DateUtilsTest.testTruncateLang59
      AssertionFailedError: expected <...MDT 2004> but was <...MST 2004>
  2x DateUtilsTest.testRound
      ArithmeticException: Calendar value too large
  1x compile error: cannot find symbol
1 candidate repeated an earlier one. No candidate passed DateUtilsTest.
\end{Verbatim}
\end{minipage}
\caption{An example handover note (illustrative values).}
\label{fig:sys-note}
\end{figure}

\paragraph{Guarantee.} Under a cold handover, with $R_2$ equal to the budget
the strong model has when it runs alone, and with the Budget Guard off,
Proposition~\ref{prop:nolost} applies: the session repairs every bug that the
strong model alone would repair on the same draws. The price is the cheap
tier's cost and delay on the bugs it fails, which
Proposition~\ref{prop:breakeven} weighs against the saving. A note gives up
this guarantee for a possible gain, and the guard gives it up for a saving.
That is why both effects are measured rather than assumed, and why a cold
handover without a stop is the default until they are.

\paragraph{Expected size.} Three figures bound the expectation. For a prompt
of $1{,}500$ tokens and candidates of $400$ tokens, a cheap call that returns
five candidates costs about one sixth of one strong call at list
prices~\citep{aibusiness2026gpt6}, so one cheap candidate costs about one
thirtieth. By the sufficient condition of Section~\ref{sec:bg-escalation}, ten
cheap candidates then pay for themselves if the cheap tier repairs more than a
third of the bugs, and five if it repairs more than a sixth. And in
neighbouring settings escalation saved between $26\,\%$ and $40\,\%$ of
cost~\citep{chen2024cascading,gandhi2025strongweak}. The expectation is
therefore a saving of tens of percent in money, not a factor. The experiments
report it together with the break-even it implies.

\paragraph{Relation to existing work.} Cascades and routers are not
new~\citep{chen2023frugalgpt,ong2024routellm,chen2024cascading,gandhi2025strongweak,son2026swerouter,madeyski2026triage}.
Three things differ here. The trigger is a verified stall rather than a
prediction, so it needs no training and cannot be wrong about a verdict. The
cheap tier's candidates are validated on a warm substrate, and that validation
is priced. And the handover is varied and measured instead of fixed.

\section{Module M3: Budget Guard}
\label{sec:sys-guard}

\paragraph{Principle.} The Budget Guard keeps the cost per bug within a budget
and reports why a session stopped. Failed sessions account for most of the
cost~\citep{fan2025sweeffi}, and current tools stop only at a fixed limit,
without giving a reason (Section~\ref{sec:rel-practice}).

\paragraph{When to stop.} Each tier has its own share of the session: $B_1$
candidates for the cheap tier and $R_2$ rounds for the strong tier. On top of
these the guard watches two things. It stops the session when the top tier
stalls, that is, after $r$ consecutive rounds of the strong tier without
progress, where progress is judged against everything the session has seen.
And it stops before a model call that the remaining budget $B$, in time,
tokens or money, cannot pay for, judged by the average cost of the tier's
calls so far. The values of $r$ and $B$ are chosen on the development set
(Section~\ref{sec:exp-prereg}).

\paragraph{Give-up report.} A session that ends without a plausible patch
returns: why it stopped; what each tier tried and what it spent; and the
failure signatures, most frequent first, with the test, the error and an
example patch. The signatures show where the models repeatedly failed, which
is the kind of information (fault location and tests) that developers find
most useful when they fix a bug themselves~\citep{noller2022trust}.
Figure~\ref{fig:sys-report} shows an example.

\begin{figure}[htbp]
\centering
\begin{minipage}{0.95\textwidth}
\begin{Verbatim}[frame=single,rulecolor=\color{boxline},fontsize=\footnotesize,framesep=2mm]
Give-up report (bug Lang-65, session 2)
Stopped: no progress for 3 rounds at the strong tier
Tiers:   cheap 10 candidates (9 distinct, 1 repeat); strong 6 rounds
Spent:   68 s (models 61 s, validation 7 s), 33k tokens, $0.07
Failures, most frequent first:
  9x DateUtilsTest.testTruncateLang59
      AssertionFailedError: expected <...MDT 2004> but was <...MST 2004>
  5x DateUtilsTest.testRound
      ArithmeticException: Calendar value too large
  1x compile error: cannot find symbol
Example patch for the most frequent failure (strong tier, round 4):
    private static void modify(Calendar val, int field, boolean round) {
        ...
    }
\end{Verbatim}
\end{minipage}
\caption{An example give-up report (illustrative values).}
\label{fig:sys-report}
\end{figure}

\paragraph{Measured apart.} Stopping is the one decision of the bot that can
lose a repair the strong model would still have made. In the experiments the
guard therefore only records when it would stop. The stop is applied
afterwards by replay, which gives its saving and its loss for every value of
$r$ and $B$ from the same sessions.

\section{When the System Helps}
\label{sec:sys-regimes}

\begin{itemize}
\item \textbf{Escalation} saves most when the price gap between the tiers is
wide and the cheap tier repairs a fair share of the bugs. If it repairs almost
none, the bot should start at the strong tier. If it repairs almost all, the
strong tier is rarely called, and the saving approaches the whole gap.
\item \textbf{The Warm Substrate} saves in proportion to the share of
validation spent building and starting. It is also what keeps cheap-first from
costing time.
\item \textbf{The Evidence Memory} saves in proportion to how often the models
repeat themselves and each other.
\item \textbf{The handover note} can help only on bugs that reach the strong
tier and are not repaired in its first round.
\item \textbf{The Budget Guard} saves most when sessions go on long after the
models are stuck.
\item \textbf{A slow strong model} is the one cost that no module removes.
Escalation only avoids calling it.
\end{itemize}

The modules therefore help in different regimes, and the regime decides their
order of importance rather than the other way round. RQ1 says which regime the
benchmark is in.
